% Geometric mean of rho = P_sinc/P_2s over 10<=kR<=40: patterns v_c, v_0..v_7 per cap
% half-angle, and the six frozen design fields at 30 degrees. Same lightness-monotone map
% as fig_regime_map_v2. Cells built in TeX from data/f8_geomean_*.dat.
\documentclass[border=1mm]{standalone}
\usepackage{amsmath,amssymb,bm}
\usepackage{jsvfigs}
\usepackage{pgfplots}
\pgfplotsset{compat=1.18}
\usepgfplotslibrary{patchplots}
\usepackage{pgfplotstable}

\definecolor{rhoa}{rgb}{0.88,0.93,0.99}
\definecolor{rhob}{rgb}{0.78,0.86,0.95}
\definecolor{rhoc}{rgb}{0.70,0.70,0.70}
\definecolor{rhod}{rgb}{0.55,0.30,0.28}
\definecolor{rhoe}{rgb}{0.15,0.04,0.04}

\pgfplotsset{
  colormap={rhomono}{color(-3)=(rhoa) color(-1.5)=(rhob) color(0)=(rhoc)
                     color(1.5)=(rhod) color(3)=(rhoe)},
}

\begin{document}
\pgfplotstableread{data/f8_geomean_modes.dat}\tabM
\pgfplotstableread{data/f8_geomean_designs30.dat}\tabD
\pgfplotstablegetrowsof{\tabM}\pgfmathtruncatemacro{\nm}{\pgfplotsretval-1}
\pgfplotstablegetrowsof{\tabD}\pgfmathtruncatemacro{\nd}{\pgfplotsretval-1}

\begin{tikzpicture}
\begin{axis}[name=axc, at={(0.9cm,1.15cm)}, anchor=south west, scale only axis,
  width=104mm, height=19mm, xmin=-0.5, xmax=15, ymin=-0.5, ymax=5.5, y dir=reverse,
  clip mode=individual, axis line style={draw=none},
  xtick={0,1,2,3,4,5,6,7,8,9.5,10.5,11.5,12.5,13.5,14.5},
  xticklabels={$v_{\rm c}$,$v_0$,$v_1$,$v_2$,$v_3$,$v_4$,$v_5$,$v_6$,$v_7$,U,B,C,D,E,G},
  xticklabel pos=upper, ytick={0,1,2,3,4,5},
  yticklabels={$15^\circ$,$30^\circ$,$60^\circ$,$90^\circ$,$120^\circ$,$150^\circ$},
  major tick length=0mm, tick style={draw=none}, every tick label/.append style={note},
  colormap name=rhomono, point meta min=-3, point meta max=3]
  \xdef\cellsC{}
  \foreach \i in {0,...,\nm}{%
    \pgfplotstablegetelem{\i}{col}\of{\tabM}\let\cx\pgfplotsretval
    \pgfplotstablegetelem{\i}{row}\of{\tabM}\let\cy\pgfplotsretval
    \pgfplotstablegetelem{\i}{l2c}\of{\tabM}\let\mm\pgfplotsretval
    \pgfmathsetmacro{\xa}{\cx-0.5}\pgfmathsetmacro{\xb}{\cx+0.5}
    \pgfmathsetmacro{\ya}{\cy-0.5}\pgfmathsetmacro{\yb}{\cy+0.5}
    \xdef\cellsC{\cellsC (\xa,\ya) [\mm] (\xb,\ya) [\mm] (\xb,\yb) [\mm] (\xa,\yb) [\mm]}
  }
  \foreach \i in {0,...,\nd}{%
    \pgfplotstablegetelem{\i}{col}\of{\tabD}\let\cx\pgfplotsretval
    \pgfplotstablegetelem{\i}{l2c}\of{\tabD}\let\mm\pgfplotsretval
    \pgfmathsetmacro{\xa}{9.5+\cx-0.5}\pgfmathsetmacro{\xb}{9.5+\cx+0.5}
    \xdef\cellsC{\cellsC (\xa,0.5) [\mm] (\xb,0.5) [\mm] (\xb,1.5) [\mm] (\xa,1.5) [\mm]}
  }
  \edef\tmp{\noexpand\addplot[patch, patch type=rectangle, shader=flat, point meta=explicit, draw=none]
     coordinates {\cellsC};}\tmp
  % printed geometric means; white digits on the dark cells
  \foreach \i in {0,...,\nm}{%
    \pgfplotstablegetelem{\i}{col}\of{\tabM}\let\cx\pgfplotsretval
    \pgfplotstablegetelem{\i}{row}\of{\tabM}\let\cy\pgfplotsretval
    \pgfplotstablegetelem{\i}{gm}\of{\tabM}\let\gv\pgfplotsretval
    \pgfplotstablegetelem{\i}{l2c}\of{\tabM}\let\mm\pgfplotsretval
    \pgfmathparse{\mm>1.2?1:0}\let\dk\pgfmathresult
    \edef\tmp{\noexpand\node[note, text=\ifnum\dk=1 white\else black\fi, anchor=center, inner sep=0pt] at (axis cs:\cx,\cy) {\gv};}\tmp
  }
  \foreach \i in {0,...,\nd}{%
    \pgfplotstablegetelem{\i}{col}\of{\tabD}\let\cx\pgfplotsretval
    \pgfplotstablegetelem{\i}{gm}\of{\tabD}\let\gv\pgfplotsretval
    \pgfplotstablegetelem{\i}{l2c}\of{\tabD}\let\mm\pgfplotsretval
    \pgfmathparse{\mm>1.2?1:0}\let\dk\pgfmathresult
    \pgfmathsetmacro{\xp}{9.5+\cx}
    \edef\tmp{\noexpand\node[note, text=\ifnum\dk=1 white\else black\fi, anchor=center, inner sep=0pt] at (axis cs:\xp,1) {\gv};}\tmp
  }
  \node[note, anchor=south, align=center] at ([yshift=4.5mm]axis cs:12,-0.5) {design fields, $\theta_0=30^\circ$};
  \node[paneltitle, anchor=south west] at ([xshift=-1.2mm,yshift=7mm]current axis.north west)
     {Geometric mean of $\rho$ over $10\le kR\le40$};
  \node[lab, rotate=90, anchor=south] at ([xshift=-7.5mm]current axis.west) {$\theta_0$};
\end{axis}

% ---------------------------------------------------------------- colour bar
\begin{scope}[shift={(4.6cm,0.5cm)}]
  \def\un{0.6}
  \shade[left color=rhoa, right color=rhob] (0,0) rectangle (1.5*\un,0.2);
  \shade[left color=rhob, right color=rhoc] (1.5*\un,0) rectangle (3*\un,0.2);
  \shade[left color=rhoc, right color=rhod] (3*\un,0) rectangle (4.5*\un,0.2);
  \shade[left color=rhod, right color=rhoe] (4.5*\un,0) rectangle (6*\un,0.2);
  \draw[witness] (0,0) rectangle (6*\un,0.2);
  \foreach \t/\s in {-3/{1/8}, -2/{1/4}, -1/{1/2}, 0/{1}, 1/{2}, 2/{4}, 3/{8}} {
    \draw[witness] ({(\t+3)*\un},0) -- ({(\t+3)*\un},-0.07);
    \node[note, anchor=north, inner sep=0.3pt] at ({(\t+3)*\un},-0.07) {\s};
  }
  \node[note, anchor=east, align=right] at (-0.15,0.1) {sinc under-predicts};
  \node[note, anchor=west, align=left] at (6*\un+0.15,0.1) {sinc over-predicts};
\end{scope}
\end{tikzpicture}
\end{document}

